\documentclass[../../main.tex]{subfiles}

\begin{document}

\section{有限集}\label{section:ch01-有限集}

或许你已经知道集合可分为有限集与无限集、可数集与不可数集等类型。然而，在本节及下一节中，我们仍要对它们加以讨论，不仅使你完全弄明白这些概念，而且还要讲一些后面将会遇到的逻辑学中的问题。首先研究有限集。

回想一下，如果 $n$ 是一个正整数，用 $S_{n}$ 表示全体小于 $n$ 的正整数的集合，并称之为正整数的一个截。集合 $S_{n}$ 就是有限集的样板。

\begin{definition}[有限集与基数]\label{definition:ch01-有限集}
集合 $A$ 称为有限的（finite），如果 $A$ 与正整数的某一个截之间存在一个一一对应。即：$A$ 称为有限的，如果 $A$ 为空集，或者对于某正整数 $n$ 存在一个一一对应
\begin{align*}
f: A \longrightarrow \{1, \cdots, n\}.
\end{align*}
前一种情形，称 $A$ 的基数（cardinality）为 $0$；后一种情形，称 $A$ 的基数为 $n$。
\end{definition}

例如，集合 $\{1, \cdots, n\}$ 的基数为 $n$，因为恒等函数是它与自身的一个一一对应。

特别要注意，我们并没有证明，对于有限集 $A$，$A$ 的基数由 $A$ 唯一决定。当然，空集的基数为 $0$。但就我们所知，可能存在集合 $A$ 与两个不同集合 $\{1, \cdots, n\}$ 和 $\{1, \cdots, m\}$ 之间的一一对应。这种可能性似乎是荒唐的，好像说两个人“数”一个盒子里的石子，得的结果不同，却又都是正确的一样。根据我们日常生活中“数”东西的经验，这是不可能的。而事实上，当 $n$ 是 $1$，$2$，$3$ 这种小的数目时，也很容易验证。但是，如果 $n$ 等于 $500$ 万，要直接证明就难以设想了。

对于如此大的 $n$，凭经验来证实是很困难的。例如，可以做这样一个试验。找一辆满载石子的货车，请十个人各自独立地去“数一数”石子的数目。就这个具体问题而言，很难想像他们“数”的结果会是一样的。当然，可以断定他们中至少有一个人“数”错了。而这其实是假定了要用经验来验证的前面那个关于基数唯一性的结果是正确的。另一种解释就是，在给定的石子集合与正整数的两个不同截之间都有一一对应。

实际生活中人们接受第一种解释。“数”东西的经验使我们相信，对于成员个数比较少的集合所产生的结论，对于任意大的集合也应该成立。

然而在数学中（与现实生活不一样）不能轻易相信这样的论断。只有用一一对应的存在性，而不是依靠具体的“数”，那才是有价值的数学证明。我们马上要证明，若 $n \neq m$，则不可能存在从给定集合 $A$ 到两个集合 $\{1, \cdots, n\}$ 与 $\{1, \cdots, m\}$ 的两个一一对应。

关于有限集，还有一些“极其显然”的事实可以用数学加以证明，本节将证明其中的一部分，而把其余的留作习题。一个简单的事情是：

\begin{lemma}\label{lemma:ch01-有限集引理}
设 $n$ 是一个正整数，$A$ 是一个集合，$a_{0}$ 是 $A$ 的一个元素。则存在集合 $A$ 与集合 $\{1, \cdots, n+1\}$ 之间的一一对应 $f$，当且仅当存在集合 $A - \{a_{0}\}$ 与 $\{1, \cdots, n\}$ 之间的一一对应 $g$。
\end{lemma}

\par
\begin{figure}[htbp]
\centering
\includegraphics[width=0.55\textwidth]{ch01-fig6-1}
\caption{引理\ref{lemma:ch01-有限集引理}的证明：交换 $a_{0}$ 与 $a_{1}$ 的像}\label{figure:ch01-有限集引理}
\end{figure}

\begin{proof}
要证明两种蕴涵关系，首先假定存在一个一一对应
\begin{align*}
g: A - \{a_{0}\} \longrightarrow \{1, \cdots, n\}.
\end{align*}
那么由下式
\begin{align*}
f(x) = g(x), \quad \text{对于 } x \in A - \{a_{0}\},
\end{align*}
\begin{align*}
f(a_{0}) = n + 1
\end{align*}
所定义的函数 $f: A \longrightarrow \{1, \cdots, n+1\}$ 就是一个一一对应。

为了证明其逆，我们假定存在一个一一对应
\begin{align*}
f: A \longrightarrow \{1, \cdots, n+1\}.
\end{align*}
如果 $f$ 恰好把 $a_{0}$ 映成数 $n+1$，事情就容易了。那时，限制映射 $f \mid A - \{a_{0}\}$ 就是 $A - \{a_{0}\}$ 与 $\{1, \cdots, n\}$ 之间的一一对应。如若不然，令 $f(a_{0}) = m$，$a_{1}$ 是 $A$ 中满足 $f(a_{1}) = n+1$ 的点，则 $a_{1} \neq a_{0}$。由
\begin{align*}
h(a_{0}) = n + 1,
\end{align*}
\begin{align*}
h(a_{1}) = m,
\end{align*}
\begin{align*}
h(x) = f(x), \quad \text{对于 } x \in A - \{a_{0}\} - \{a_{1}\}
\end{align*}
定义一个新函数
\begin{align*}
h: A \longrightarrow \{1, \cdots, n+1\}.
\end{align*}
容易验证 $h$ 是一个一一对应（参见\cref{figure:ch01-有限集引理}）。


这就回到了简单的情形，限制映射 $h \mid A - \{a_{0}\}$ 就是 $A - \{a_{0}\}$ 与 $\{1, \cdots, n\}$ 之间的一一对应。
\end{proof}

由这个引理，可得到几个有用的推论。

\begin{theorem}[有限集的真子集]\label{theorem:ch01-有限真子集}
设 $A$ 是一个集合。假定对于某一个 $n \in \mathbb{Z}_{+}$，存在一个一一对应 $f: A \longrightarrow \{1, \cdots, n\}$。设 $B$ 为 $A$ 的一个真子集，则不存在一一对应 $g: B \longrightarrow \{1, \cdots, n\}$。但是（假定 $B \neq \varnothing$），必定有一个一一对应 $h: B \longrightarrow \{1, \cdots, m\}$ 对于某一个 $m < n$ 成立。
\end{theorem}

\begin{proof}
当 $B = \varnothing$ 时，由于空集 $B$ 与非空集 $\{1, \cdots, n\}$ 之间不存在一一对应，所以此时结论显然成立。

我们用“归纳法”来证明定理。设 $C$ 为 $\mathbb{Z}_{+}$ 的子集，它由使定理成立的那些整数 $n$ 组成。下面证明 $C$ 是包含 $1$ 的归纳集。由此可见 $C = \mathbb{Z}_{+}$，从而定理对于所有正整数 $n$ 成立。

首先证明 $n = 1$ 时定理成立。这时 $A$ 仅由一个元素 $\{a\}$ 组成，它的唯一真子集 $B$ 是空集。

假定对于 $n$ 定理成立，证明定理对于 $n+1$ 成立。设 $f: A \longrightarrow \{1, \cdots, n+1\}$ 是一个一一对应，$B$ 是 $A$ 的一个非空真子集。在 $B$ 与 $A - B$ 中各取一点 $a_{0}$ 与 $a_{1}$。根据前面一个引理可见，存在一个一一对应
\begin{align*}
g: A - \{a_{0}\} \longrightarrow \{1, \cdots, n\}.
\end{align*}
因为 $a_{1}$ 属于 $A - \{a_{0}\}$ 而不属于 $B - \{a_{0}\}$，所以 $B - \{a_{0}\}$ 是 $A - \{a_{0}\}$ 的一个真子集。而已经假定了定理对于整数 $n$ 成立。所以
\begin{enumerate}[(1)]
\item 不存在一一对应 $h: B - \{a_{0}\} \longrightarrow \{1, \cdots, n\}$.
\item 或者 $B - \{a_{0}\} = \varnothing$，或者对某一个 $p < n$ 存在一个一一对应
\begin{align*}
k: B - \{a_{0}\} \longrightarrow \{1, \cdots, p\}.
\end{align*}
\end{enumerate}
根据前面的引理及 (1) 可得，不存在 $B$ 与 $\{1, \cdots, n+1\}$ 之间的一一对应。这就是我们要证明的前半段。为了证明后半段，注意到如果 $B - \{a_{0}\} = \varnothing$，则存在 $B$ 与集合 $\{1\}$ 之间的一一对应，如果 $B - \{a_{0}\} \neq \varnothing$，应用前面的引理及 (2) 可见，存在 $B$ 与 $\{1, \cdots, p+1\}$ 之间的一个一一对应。两种情况都表明，对于某一个 $m < n+1$，在 $B$ 与 $\{1, \cdots, m\}$ 之间有一个一一对应。由归纳原理就证明了定理对于所有 $n \in \mathbb{Z}_{+}$ 成立。
\end{proof}

\begin{corollary}\label{corollary:ch01-有限不自映射}
如果 $A$ 是一个有限集，则不存在 $A$ 与其真子集之间的一一对应。
\end{corollary}

\begin{proof}
设 $B$ 是 $A$ 的一个真子集，$f: A \longrightarrow B$ 是一个一一对应，则根据假定，对于某一个 $n$，存在一个一一对应 $g: A \longrightarrow \{1, \cdots, n\}$。它们的复合 $g \circ f^{-1}$ 是 $B$ 与 $\{1, \cdots, n\}$ 之间的一个一一对应。与上述\cref{theorem:ch01-有限真子集}矛盾。
\end{proof}

\begin{corollary}\label{corollary:ch01-Z无限}
$\mathbb{Z}_{+}$ 不是有限集。
\end{corollary}

\begin{proof}
通过 $f(n) = n+1$ 所定义的函数 $f: \mathbb{Z}_{+} \longrightarrow \mathbb{Z}_{+} - \{1\}$ 便是 $\mathbb{Z}_{+}$ 与其一个真子集之间的一个一一对应。
\end{proof}

\begin{corollary}\label{corollary:ch01-基数唯一}
有限集 $A$ 的基数由 $A$ 唯一决定。
\end{corollary}

\begin{proof}
证明设 $m < n$。假若存在一一对应
\begin{align*}
f: A \longrightarrow \{1, \cdots, n\},
\end{align*}
\begin{align*}
g: A \longrightarrow \{1, \cdots, m\}.
\end{align*}
则其复合
\begin{align*}
g \circ f^{-1}: \{1, \cdots, n\} \longrightarrow \{1, \cdots, m\}
\end{align*}
是有限集 $\{1, \cdots, n\}$ 与其一个真子集之间的一一对应，与前面的\cref{corollary:ch01-有限不自映射}矛盾。
\end{proof}

\begin{corollary}\label{corollary:ch01-有限子集}
如果 $B$ 是有限集 $A$ 的一个子集，则 $B$ 是一个有限集。如果 $B$ 是 $A$ 的一个真子集，则 $B$ 的基数小于 $A$ 的基数。
\end{corollary}

\begin{corollary}\label{corollary:ch01-有限等价条件}
设 $B$ 是一个非空集。则下列条件等价：
\begin{enumerate}[(1)]
\item $B$ 是一个有限集.
\item 存在从正整数的某一个截到 $B$ 上的满射.
\item 存在从 $B$ 到正整数的某一个截的单射.
\end{enumerate}
\end{corollary}

\begin{proof}
(1) $\Longrightarrow$ (2)。由于 $B$ 非空，对于某个 $n$，存在一一对应 $f: \{1, \cdots, n\} \longrightarrow B$。

(2) $\Longrightarrow$ (3)。若 $f: \{1, \cdots, n\} \longrightarrow B$ 为满射，用
\begin{align*}
g(b) = f^{-1}(\{b\}) \text{ 的最小元}
\end{align*}
定义一个函数 $g: B \longrightarrow \{1, \cdots, n\}$。由于 $f$ 是一个满射，集合 $f^{-1}(\{b\})$ 非空。根据 $\mathbb{Z}_{+}$ 的良序性质可见，$g(b)$ 是唯一确定的。由于当 $b \neq b'$ 时，集合 $f^{-1}(\{b\})$ 与 $f^{-1}(\{b'\})$ 无交，因此它们的最小元一定不同，于是 $g$ 是一个单射。

(3) $\Longrightarrow$ (1)。若 $g: B \longrightarrow \{1, \cdots, n\}$ 是一个单射，则改变 $g$ 的值域可以得到从 $B$ 到 $\{1, \cdots, n\}$ 的某子集之间的一个一一对应。根据前面的推论，可见 $B$ 是一个有限集。
\end{proof}

\begin{corollary}\label{corollary:ch01-有限并与积}
若干个有限集的有限并及有限笛卡儿积是有限集。
\end{corollary}

\begin{proof}
首先证明：若 $A$ 和 $B$ 都是有限集，则 $A \cup B$ 是一个有限集。当 $A$ 和 $B$ 至少有一个是空集，结论显然成立。当 $A$ 和 $B$ 都不是空集时，对于某 $m$ 和 $n$，存在一一对应 $f: \{1, \cdots, m\} \longrightarrow A$ 及一一对应 $g: \{1, \cdots, n\} \longrightarrow B$。用
\begin{align*}
h(i) = f(i), i = 1, 2, \cdots, m \quad \text{和}
\end{align*}
\begin{align*}
h(i) = g(i - m), i = m + 1, \cdots, m + n
\end{align*}
定义一个函数 $h: \{1, \cdots, m+n\} \longrightarrow A \cup B$。易见，$h$ 是一个满射，所以 $A \cup B$ 是一个有限集。

其次，用归纳法证明：$A_{1}, \cdots, A_{n}$ 各个集合的有限性蕴涵着它们的并的有限性。$n = 1$ 时显然成立。假定对于 $n-1$ 命题成立。注意到 $A_{1} \cup \cdots \cup A_{n}$ 可以表示成两个有限集 $A_{1} \cup \cdots \cup A_{n-1}$ 与 $A_{n}$ 之并，由上段的讨论，命题成立。

再次，证明两个有限集 $A$ 和 $B$ 的笛卡儿积是有限集。任意选取 $a \in A$，由于 $\{a\} \times B$ 的点与 $B$ 的点有一一对应关系，所以 $\{a\} \times B$ 为有限集。而 $A \times B$ 可以表示成具有 $\{a\} \times B$ 形式的集合之并，并且此种形式的集合仅有有限个，从而 $A \times B$ 为有限个有限集之并，因而是有限集。

最后，用归纳法易证：若每一个 $A_{i}$ 有限，则笛卡儿积 $A_{1} \times \cdots \times A_{n}$ 也是有限集。
\end{proof}

\subsection*{习题}

\begin{exercise}\label{exercise:ch01-有限集-1}
(a) 写出所有单射
\begin{align*}
f: \{1, 2, 3\} \longrightarrow \{1, 2, 3, 4\},
\end{align*}
证明它们都不是一一对应。（这里给出了基数为 $3$ 的集合 $A$ 其基数不为 $4$ 的一个直接的证明。）

(b) 有多少个单射
\begin{align*}
f: \{1, \cdots, 8\} \longrightarrow \{1, \cdots, 10\}?
\end{align*}
（你会发现我们为什么不试图直接证明在这两个集合之间不存在一一对应。）
\end{exercise}

\begin{exercise}\label{exercise:ch01-有限集-2}
证明：若 $B$ 不是有限集，并且 $B \subset A$，则 $A$ 也不是有限集。
\end{exercise}

\begin{exercise}\label{exercise:ch01-有限集-3}
设 $X$ 为二元素集 $\{0, 1\}$。求出 $X^{\omega}$ 与其一个真子集之间的一一对应。
\end{exercise}

\begin{exercise}\label{exercise:ch01-有限集-4}
设 $A$ 是一个非空的有限全序集。
\begin{enumerate}[(a)]
\item 证明 $A$ 有一个最大元. [提示：对 $A$ 的基数进行归纳.]
\item 证明 $A$ 具有正整数的一个截的序型.
\end{enumerate}
\end{exercise}

\begin{exercise}\label{exercise:ch01-有限集-5}
$A \times B$ 为有限集是否蕴涵着 $A$ 与 $B$ 都是有限集？
\end{exercise}

\begin{exercise}\label{exercise:ch01-有限集-6}
(a) 设 $A = \{1, \cdots, n\}$。证明：在 $\mathcal{P}(A)$ 与笛卡儿积 $X^{n}$ 之间存在一个一一对应，其中 $X = \{0, 1\}$ 为二元素集。

(b) 证明：若 $A$ 是一个有限集，则 $\mathcal{P}(A)$ 也是有限集。
\end{exercise}

\begin{exercise}\label{exercise:ch01-有限集-7}
若 $A$ 和 $B$ 都是有限的，证明所有函数 $f: A \longrightarrow B$ 的集合是一个有限集。
\end{exercise}

\end{document}